\phantomsection
\part{写像}

\begin{abstract}
この章では，集合とならんで現代数学の基礎をなす「写像」を導入する．まず写像を「各元にただひとつの元を対応させる規則」として定義し，自動販売機などの身近な例で直感的な理解を深める．続いて，直積と記号$\exists!$を用いた写像の厳密な定義を紹介する．次に，像・値域・逆像といった写像に付随する集合を定義する．そのうえで，単射・全射・全単射という写像の重要な性質を，図と具体例をもとに確認し，最後に像・逆像と集合演算の関係を整理しよう．
\end{abstract}

\phantomsection
\section{写像の定義と例}

% --- この章の写像図で共通して使う TikZ スタイル ---
\tikzset{
  mapfig/.style     ={>=Latex, font=\small, line cap=round, line join=round},
  mapset/.style     ={draw=black!60, fill=black!4, line width=0.7pt},
  mapsub/.style     ={draw=black!50, fill=black!14, line width=0.5pt},
  mapelem/.style    ={circle, fill=black!85, inner sep=1.3pt},
  maparrow/.style   ={-{Latex[length=1.8mm,width=1.3mm]}, draw=black!70,
                      line width=0.55pt, shorten <=2.5pt, shorten >=3pt},
  mapleader/.style  ={draw=black!50, line width=0.4pt},
  mapsetname/.style ={font=\small, text=black!70},
  maplabL/.style    ={anchor=east, xshift=-2.5pt},
  maplabR/.style    ={anchor=west, xshift=2.5pt},
}

\phantomsection
\subsection{写像の定義}

いままでの数学の学習の過程で「関数\index{かんすう@関数}」という言葉に出会ったことがある読者もいるかもしれない．
「関数」と言われると，$y = \sin x$や$y = x^3 + 1$のような「$x$と$y$を結ぶ式」，あるいはそのグラフを思い浮かべる読者が多いであろう．
しかし，これから見るように，関数の本質は式やグラフそのものではなく，「元にただひとつの元を対応させる規則」である．
式やグラフは，この規則を表現し，視覚化するための手段にすぎない．

関数の一般的な表現として「写像」がある\footnote{現代数学では，「関数」と「写像」はほぼ同義に用いられることも多い．本書では，一般の集合の間の対応であることを強調するときに「写像」という語を用いる．なお，定義域や終域が$\mathbb{R}$など「数」の集合である写像を「関数」とよぶことが多いが，これは絶対的な区別ではない．}．写像の厳密な定義はのちほど（本章「写像の厳密な定義」の項で）述べるとして，まず簡潔な定義を提示しよう．

\begin{definition}{写像}{写像}
  集合$A$から集合$B$への\emph{写像}\index{しゃぞう@写像}$f$とは，任意の$a \in A$に対して，$b \in B$をただひとつ対応させる規則のことである．

  このとき，$A$を$f$の\emph{定義域}\index{ていぎいき@定義域}（domain），$B$を$f$の\emph{終域}\index{しゅういき@終域}（codomain）とよぶ\footnote{定義域を「始集合」，終域を「終集合」とよぶ文献もある．}．

  これを次のように表す：
  \[
    f \colon A \to B
  \]
  この写像$f$によって，$a \in A$が$b \in B$に対応するとき，
  \[
    b=f(a)
  \]
  または
  \[
    f \colon a \mapsto b
  \]
  とかく\footnote{これを$ f \colon A \ni a \mapsto b \in B$とかくときもある．}．
\end{definition}

\begin{wrapfigure}{r}{6.0cm}
  \centering
  \begin{tikzpicture}[mapfig]
    % 集合
    \draw[mapset] (-1.75,0) ellipse [x radius=0.95, y radius=1.25];
    \draw[mapset] ( 1.75,0) ellipse [x radius=0.95, y radius=1.25];
    \node[mapsetname] at (-1.75,1.55) {$A$};
    \node[mapsetname] at ( 1.75,1.55) {$B$};
    \node[mapsetname] at (0,1.35) {$f$};
    % 要素
    \coordinate (A1) at (-1.75, 0.75);
    \coordinate (A2) at (-1.75, 0);
    \coordinate (A3) at (-1.75,-0.75);
    \coordinate (Ba) at ( 1.75, 0.75);
    \coordinate (Bb) at ( 1.75, 0);
    \coordinate (Bc) at ( 1.75,-0.75);
    \foreach \p in {A1,A2,A3,Ba,Bb,Bc} \node[mapelem] at (\p) {};
    \node[maplabL] at (A1) {$1$};
    \node[maplabL] at (A2) {$2$};
    \node[maplabL] at (A3) {$3$};
    \node[maplabR] at (Ba) {$a$};
    \node[maplabR] at (Bb) {$b$};
    \node[maplabR] at (Bc) {$c$};
    % 対応
    \draw[maparrow] (A1) to[bend left=10] (Ba);
    \draw[maparrow] (A2) -- (Bb);
    \draw[maparrow] (A3) to[bend left=12] (Ba);
  \end{tikzpicture}
  \caption{写像 $f \colon A \to B$}
\end{wrapfigure}

$A=\{ 1,2 ,3 \}$， $B=\{ a, b,c \}$とする．$f(1) = a$，$f(2) = b$，$f(3) = a$とすると，この写像$f$は右図のように表現できる．


この例では，$1$と$3$がともに$a$に対応しているが，上記の写像の定義ではこのような場合があってもよい．
また，$A$のどの元にも対応しない元$c$の存在も許容される．

学習が進むと，「任意の$ a_1 , a_2 \in A$に対して，$ a_1 \ne a_2$ならば$f(a_1) \ne f(a_2)$である写像」や
「任意の$ b \in B$に対して，ある$a \in A$が存在して，$b=f(a)$となる写像」を考える場合もあるが，この説明はもう少しあとにしよう．

\phantomsection
\subsection{写像の例}

ここでは直感的な理解を深めるために，具体例を挙げてみよう．

\begin{example}{自動販売機}{自動販売機}
  自動販売機のボタンの集合を$A$，提供される飲み物の集合を$B$とする．
  各ボタン$a \in A$に，そのボタンを押すと出てくる飲み物$f(a) \in B$を対応させる．
  どのボタンにも飲み物がただひとつ定まるから，$f \colon A \to B$は写像である．
\end{example}

\begin{example}{写像の例と写像でない例}{写像の例と写像でない例}
  \begin{itemize}
    \item $f \colon \mathbb{R} \ni x \mapsto x^2 \in \mathbb{R}$は写像である．
          各実数$x$に対して，$x^2$がただひとつ定まるからである．
    \item 各正の実数$x$に「2乗すると$x$になる実数」を対応させる規則は，写像ではない．
          たとえば$x = 4$には$2$と$-2$の2つが対応し，「ただひとつ対応させる」という条件を満たさないからである．
  \end{itemize}
\end{example}

\begin{remark}[「式」と「写像」の違い]
  「多項式$x^2 + x$は関数である」という言い方をすることがあるが，正確には，式$x^2 + x$それ自体は単なる多項式であって，写像ではない．定義域と終域を定め，
  \[
    f \colon \mathbb{R} \ni x \mapsto x^2 + x \in \mathbb{R}
  \]
  のように対応の規則を与えてはじめて写像となる．同じ式を用いても，定義域を$[0, 1]$に変えれば，写像としては別のものである．
  逆に，\exref{自動販売機}のように，ひとつの式では表されない写像も存在する．すなわち，「写像であること」と「式で表されること」は別のことがらである．

  また，関数をグラフでとらえる見方も有用であるが，グラフ（点$(x, f(x))$全体のなす集合）だけでは終域の情報は定まらないことに注意したい．
  同じ式・同じグラフでも，終域の取り方によって写像の性質は変わりうる（この違いは，\deref{全射}のあとの例で実際に現れる）．
\end{remark}


\phantomsection
\subsection{写像の厳密な定義}

\deref{写像}では，写像を「元にただひとつの元を対応させる\emph{規則}」として定義した．この定義は直感的でわかりやすいが，「規則」という言葉自体は数学的に定義されていない，という弱点がある．そこで本項では，集合の言葉だけを用いた写像の厳密な定義を紹介する．準備として，順序対・直積と，記号$\exists!$を導入しよう．

\phantomsection
\begin{definition}{順序対と直積}{順序対と直積}
  $a \in A$と$b \in B$を，この順序を込めて組にしたもの$(a, b)$を\emph{順序対}\index{じゅんじょつい@順序対}という．順序対については，
  \[
    (a_1, b_1) = (a_2, b_2)
    \lequiv
    (a_1 = a_2) \land (b_1 = b_2)
  \]
  が成り立つものとする\footnote{順序対そのものも，$(a, b) \coloneqq \{ \{a\},\, \{a, b\} \}$のように集合として定義できる（クラトフスキーの定義）．本書では，上の性質を満たすものとして順序対を扱えば十分である．}．順序対全体からなる集合
  \[
    A \times B \coloneqq \{ (a, b) \mid a \in A,\ b \in B \}
  \]
  を$A$と$B$の\emph{直積}\index{ちょくせき@直積}という．
\end{definition}

集合においては$\{1, 2\} = \{2, 1\}$と要素の順序を問わなかったのに対し，順序対では$(1, 2) \ne (2, 1)$である．座標平面の点$(x, y)$は順序対の身近な例であり，実際，座標平面は直積$\mathbb{R} \times \mathbb{R}$とみなすことができる．

\phantomsection
\begin{definition}{記号$\exists!$}{一意存在の記号}
  $P(y)$を述語とする．「$P(y)$を満たす$y \in B$が\textbf{ただひとつ}存在する」ことを
  \[
    \lexu{y \in B} \lst P(y)
  \]
  と表す\index{いちいそんざい@一意存在}．これは，
  \[
    \lex{y \in B} \lst \Bigl( P(y) \land \bigl( \lall{z \in B} \lst ( P(z) \limp z = y ) \bigr) \Bigr)
  \]
  の略記である．
\end{definition}

すなわち，$\lexu{y \in B} \lst P(y)$は，
\begin{itemize}
  \item \textbf{存在}：$P(y)$を満たす$y \in B$が（少なくともひとつ）存在する
  \item \textbf{一意性}：$P(z)$を満たす$z \in B$は，その$y$のほかにない
\end{itemize}
というふたつの主張をあわせたものである．\partref{証明}で見たように，「存在」と「一意性」は別々に証明すべきふたつの課題であった．記号$\exists!$は，まさにその両方を一度に主張する記号なのである．

準備が整ったので，写像の厳密な定義を述べよう．

\phantomsection
\begin{definition}{写像（厳密な定義）}{写像の厳密な定義}
  $A$，$B$を集合とする．$A$から$B$への\emph{写像}$f$とは，直積$A \times B$の部分集合$G \subset A \times B$であって，
  \[
    \lall{a \in A} \lsep \lexu{b \in B} \lst (a, b) \in G
  \]
  を満たすもののことである．このとき，各$a \in A$に対して，$(a, b) \in G$を満たすただひとつの$b \in B$を$f(a)$とかく．
\end{definition}

\begin{remark}[ふたつの定義の関係]
  \deref{写像}の「任意の$a \in A$に対して，$b \in B$をただひとつ対応させる」という条件は，厳密な定義では論理式
  \[
    \lall{a \in A} \lsep \lexu{b \in B} \lst (a, b) \in G
  \]
  として表現されている．そして，「規則」という無定義の言葉が，「対応する組$(a, b)$をすべて集めた集合$G$」に置き換えられている点が重要である．たとえば，$f \colon \mathbb{R} \ni x \mapsto x^2 \in \mathbb{R}$に対応する集合は
  \[
    G = \{ (x, x^2) \mid x \in \mathbb{R} \} \subset \mathbb{R} \times \mathbb{R}
  \]
  であり，これはまさに$y = x^2$の\emph{グラフ}である．すなわち，厳密な定義は「写像とは，そのグラフのことである」という見方を与えている\footnote{ただし，集合$G$だけからは終域$B$は復元できない（「式」と「写像」の違いに関する先の注意を参照）．そこで，終域の情報まで込めて，写像を3つ組$(A, B, G)$として定義する流儀もある．}．
\end{remark}

なお，本書では以後も\deref{写像}の直感的な定義に基づいて議論を進めるが，「規則」という言葉が集合の言葉で正当化できることは，記憶にとどめておいてほしい．

\subsection{像と値域，逆像}

\begin{definition}{像と値域}{像と値域}
  写像$f \colon A \to B$において，集合$A$の部分集合$A' \subset A$に対して，
  \[
    f(A') = \{ f(a) \mid a \in A' \}
  \]
  で定義される集合を，$f$による$A'$の\emph{像}という．

  特に，定義域$A$全体の像
  \[
    f(A) = \{ f(a) \mid a \in A \}
  \]
  を，写像$f$の\emph{値域}という\footnote{「値域」という語は文献によって流儀が異なり，終域$B$の意味で用いられることもある．本書では，$f(A)$を値域とよぶ．}．

  \index{ぞう@像} \index{ちいき@値域}
\end{definition}


\begin{definition}{逆像}{逆像}
  写像$f \colon A \to B$と集合$B' \subset B$に対して，$B'$の\emph{逆像}とは，
  \[
    f^{-1}(B') = \{ a \in A \mid f(a) \in B' \}
  \]
  で定義される集合である．

  \index{ぎゃくぞう@逆像}
\end{definition}

\begin{remark}[逆像と逆写像]
  \deref{逆像}で定義した逆像は「写像」ではなく，あくまで「集合」であることに注意したい．
  逆像と紛らわしい単語に「逆写像」があるが，逆写像は$f$が全単射であるときに定義される写像である．その一方で，逆像は$f$が全単射でなくても定義される．
\end{remark}

\begin{figure}[ht]
  \centering
  \begin{minipage}[t]{0.48\linewidth}
    \centering
    \begin{tikzpicture}[mapfig]
      % 集合
      \draw[mapset] (-1.75,0) ellipse [x radius=0.95, y radius=1.25];
      \draw[mapset] ( 1.75,0) ellipse [x radius=0.95, y radius=1.25];
      \node[mapsetname] at (-1.75,-1.55) {$A$};
      \node[mapsetname] at ( 1.75,-1.55) {$B$};
      \node[mapsetname] at (0,1.35) {$f$};
      % A' = {1, 2} とその像 f(A') = {a, b}
      \draw[mapsub] (-1.75,0.375) ellipse [x radius=0.55, y radius=0.85];
      \draw[mapsub] ( 1.75,0.125) ellipse [x radius=0.50, y radius=0.85];
      \draw[mapleader] (-2.14,0.98) -- (-2.45,1.35)
        node[anchor=south, font=\small] {$A'$};
      \draw[mapleader] ( 2.10,0.72) -- ( 2.45,1.35)
        node[anchor=south, font=\small] {$f(A')$};
      % 要素
      \coordinate (A1) at (-1.75, 0.75);
      \coordinate (A2) at (-1.75, 0);
      \coordinate (A3) at (-1.75,-0.75);
      \coordinate (Ba) at ( 1.75, 0.5);
      \coordinate (Bb) at ( 1.75,-0.25);
      \foreach \p in {A1,A2,A3,Ba,Bb} \node[mapelem] at (\p) {};
      \node[maplabL] at (A1) {$1$};
      \node[maplabL] at (A2) {$2$};
      \node[maplabL] at (A3) {$3$};
      \node[maplabR] at (Ba) {$a$};
      \node[maplabR] at (Bb) {$b$};
      % 対応
      \draw[maparrow] (A1) to[bend left=10] (Ba);
      \draw[maparrow] (A2) -- (Bb);
      \draw[maparrow] (A3) to[bend left=12] (Ba);
    \end{tikzpicture}
    \caption{$A' = \{1, 2\}$ の像 $f(A') = \{a, b\}$}
  \end{minipage}
  \hfill
  \begin{minipage}[t]{0.48\linewidth}
    \centering
    \begin{tikzpicture}[mapfig]
      % 集合
      \draw[mapset] (-1.75,0) ellipse [x radius=0.95, y radius=1.35];
      \draw[mapset] ( 1.75,0) ellipse [x radius=0.95, y radius=1.35];
      \node[mapsetname] at (-1.75,-1.65) {$A$};
      \node[mapsetname] at ( 1.75,-1.65) {$B$};
      \node[mapsetname] at (0,1.35) {$f$};
      % B' = {a} とその逆像 f^{-1}(B') = {1, 2, 3}
      \draw[mapsub] ( 1.75,0.5) ellipse [x radius=0.45, y radius=0.42];
      \draw[mapsub] (-1.75,0.30) ellipse [x radius=0.60, y radius=1.00];
      \draw[mapleader] (-2.17,1.01) -- (-2.45,1.45)
        node[anchor=south, font=\small] {$f^{-1}(B')$};
      \draw[mapleader] ( 2.07,0.80) -- ( 2.40,1.45)
        node[anchor=south, font=\small] {$B'$};
      % 要素
      \coordinate (A1) at (-1.75, 0.95);
      \coordinate (A2) at (-1.75, 0.30);
      \coordinate (A3) at (-1.75,-0.35);
      \coordinate (A4) at (-1.75,-1.00);
      \coordinate (Ba) at ( 1.75, 0.5);
      \coordinate (Bb) at ( 1.75,-0.5);
      \foreach \p in {A1,A2,A3,A4,Ba,Bb} \node[mapelem] at (\p) {};
      \node[maplabL] at (A1) {$1$};
      \node[maplabL] at (A2) {$2$};
      \node[maplabL] at (A3) {$3$};
      \node[maplabL] at (A4) {$4$};
      \node[maplabR] at (Ba) {$a$};
      \node[maplabR] at (Bb) {$b$};
      % 対応
      \draw[maparrow] (A1) to[bend left=10] (Ba);
      \draw[maparrow] (A2) -- (Ba);
      \draw[maparrow] (A3) to[bend right=8] (Ba);
      \draw[maparrow] (A4) to[bend right=12] (Bb);
    \end{tikzpicture}
    \caption{$B' = \{a\}$ の逆像 $f^{-1}(B') = \{1, 2, 3\}$}
  \end{minipage}
\end{figure}



\subsection{単射\index{たんしゃ@単射}}

\begin{definition}{単射}{単射}
  写像$f \colon A \to B$が\emph{単射}であるとは，任意の$a_1, a_2 \in A$に対して，$a_1 \ne a_2$ならば$f(a_1) \ne f(a_2)$が成り立つことである\footnote{この定義は，対偶を考えると「$f(a_1) =f(a_2)$ならば$a_1=a_2$が成り立つ」とも言い換えられる．}．
\end{definition}

\phantomsection
\begin{example}[
    sidebyside,
    lefthand width=7.4cm,
    righthand width=6.2cm,
    sidebyside gap=4mm,
  ]{単射の図解}{単射の図解}
  集合$A = \{1, 2, 3\}$，$B = \{a, b, c, d\}$と写像
  \[
    f(1) = a, \quad f(2) = b, \quad f(3) = c
  \]
  を考える．
  $A$の異なる元は$B$の異なる元にうつるから，$f$は単射である．
  しかし，この$f$は全射でない．

  \tcblower

  \centering
  \begin{tikzpicture}[mapfig]
    % 集合
    \draw[mapset] (-1.75,0) ellipse [x radius=0.95, y radius=1.35];
    \draw[mapset] ( 1.75,0) ellipse [x radius=0.95, y radius=1.35];
    \node[mapsetname] at (-1.75,1.65) {$A$};
    \node[mapsetname] at ( 1.75,1.65) {$B$};
    \node[mapsetname] at (0,1.4) {$f$};
    % 要素
    \coordinate (A1) at (-1.75, 0.75);
    \coordinate (A2) at (-1.75, 0);
    \coordinate (A3) at (-1.75,-0.75);
    \coordinate (Ba) at ( 1.75, 0.95);
    \coordinate (Bb) at ( 1.75, 0.32);
    \coordinate (Bc) at ( 1.75,-0.32);
    \coordinate (Bd) at ( 1.75,-0.95);
    \foreach \p in {A1,A2,A3,Ba,Bb,Bc,Bd} \node[mapelem] at (\p) {};
    \node[maplabL] at (A1) {$1$};
    \node[maplabL] at (A2) {$2$};
    \node[maplabL] at (A3) {$3$};
    \node[maplabR] at (Ba) {$a$};
    \node[maplabR] at (Bb) {$b$};
    \node[maplabR] at (Bc) {$c$};
    \node[maplabR] at (Bd) {$d$};
    % 対応
    \draw[maparrow] (A1) to[bend left=8] (Ba);
    \draw[maparrow] (A2) to[bend left=6] (Bb);
    \draw[maparrow] (A3) to[bend right=6] (Bc);
  \end{tikzpicture}
\end{example}

\begin{example}{単射の例}{単射の例}
  \[
    f \colon \mathbb{N} \ni n \mapsto 2n \in \mathbb{N}
  \]
  は単射である．実際，$2n_1 = 2n_2$ならば$n_1 = n_2$である．
\end{example}

\begin{example}{電話番号と電話機}{電話番号と電話機}
  電話番号の集合を$A$，電話機の集合を$B$とし，各電話番号$a \in A$に対応する電話機$f(a) \in B$を対応させる写像$f \colon A \to B$を考える．
  異なる電話番号が同じ電話機を指すことはないから，$f$は単射である．
\end{example}

\subsection{全射\index{ぜんしゃ@全射}}

\begin{definition}{全射}{全射}
  写像$f \colon A \to B$が\emph{全射}であるとは，任意の$b \in B$に対して，ある$a \in A$が存在して，$f(a) = b$が成り立つことである．
\end{definition}

\phantomsection
\begin{example}[
    sidebyside,
    lefthand width=7.4cm,
    righthand width=6.2cm,
    sidebyside gap=4mm,
  ]{全射の図解}{全射の図解}
  集合$A = \{1, 2, 3\}$，$B = \{a, b\}$と写像
  \[
    f(1) = a, \quad f(2) = a, \quad f(3) = b
  \]
  を考える．
  $B$のどの元も$A$のある元の像になっているから，$f$は全射である．
  しかし，$f(1) = f(2) = a$であるから，$f$は単射でない．

  \tcblower

  \centering
  \begin{tikzpicture}[mapfig]
    % 集合
    \draw[mapset] (-1.75,0) ellipse [x radius=0.95, y radius=1.25];
    \draw[mapset] ( 1.75,0) ellipse [x radius=0.95, y radius=1.25];
    \node[mapsetname] at (-1.75,1.55) {$A$};
    \node[mapsetname] at ( 1.75,1.55) {$B$};
    \node[mapsetname] at (0,1.35) {$f$};
    % 要素
    \coordinate (A1) at (-1.75, 0.75);
    \coordinate (A2) at (-1.75, 0);
    \coordinate (A3) at (-1.75,-0.75);
    \coordinate (Ba) at ( 1.75, 0.5);
    \coordinate (Bb) at ( 1.75,-0.5);
    \foreach \p in {A1,A2,A3,Ba,Bb} \node[mapelem] at (\p) {};
    \node[maplabL] at (A1) {$1$};
    \node[maplabL] at (A2) {$2$};
    \node[maplabL] at (A3) {$3$};
    \node[maplabR] at (Ba) {$a$};
    \node[maplabR] at (Bb) {$b$};
    % 対応
    \draw[maparrow] (A1) to[bend left=10] (Ba);
    \draw[maparrow] (A2) -- (Ba);
    \draw[maparrow] (A3) to[bend right=8] (Bb);
  \end{tikzpicture}
\end{example}

\begin{example}{全射の例}{全射の例}
  \[
    f \colon \mathbb{R} \ni x \mapsto x^2 \in [0, \infty)
  \]
  は全射である．実際，任意の$b \in [0, \infty)$に対して，$x = \sqrt{b}$とおけば$f(x) = b$である．
  一方，$f(-1) = f(1) = 1$であるから，$f$は単射でない．

  なお，\exref{写像の例と写像でない例}の$x \mapsto x^2$とは式が同じで終域だけが異なるが，終域を$\mathbb{R}$とすると全射でなくなる．写像の性質は，式だけでなく定義域・終域を込めて定まるのである．
\end{example}

\begin{example}{分類コードと本}{分類コードと本}
  図書館にある本の集合を$A$，使用されている分類コードの集合を$B$とし，各本$a \in A$にその分類コード$f(a) \in B$を対応させる写像$f \colon A \to B$を考える．
  $B$は使用されているコード全体であるから，どのコード$b \in B$に対しても，それが割り当てられた本$a \in A$が存在する．よって$f$は全射である．
  一方，同じ分類コードを持つ本は複数ありうるから，$f$は一般に単射でない．
\end{example}


\subsection{全単射\index{ぜんたんしゃ@全単射}}

\begin{definition}{全単射}{全単射}
  写像$f \colon A \to B$が\emph{全単射}であるとは，$f$が単射かつ全射となることである．
\end{definition}

\phantomsection
\begin{example}[
    sidebyside,
    lefthand width=7.4cm,
    righthand width=6.2cm,
    sidebyside gap=4mm,
  ]{全単射の図解}{全単射の図解}
  集合$A = \{1, 2, 3\}$，$B = \{a, b, c\}$と写像
  \[
    f(1) = c, \quad f(2) = b, \quad f(3) = a
  \]
  を考える．
  この$f$は単射かつ全射，すなわち全単射である．
  $A$の元と$B$の元が，$f$によって過不足なく1対1に対応している．

  \tcblower

  \centering
  \begin{tikzpicture}[mapfig]
    % 集合
    \draw[mapset] (-1.75,0) ellipse [x radius=0.95, y radius=1.25];
    \draw[mapset] ( 1.75,0) ellipse [x radius=0.95, y radius=1.25];
    \node[mapsetname] at (-1.75,1.55) {$A$};
    \node[mapsetname] at ( 1.75,1.55) {$B$};
    \node[mapsetname] at (0,1.35) {$f$};
    % 要素
    \coordinate (A1) at (-1.75, 0.75);
    \coordinate (A2) at (-1.75, 0);
    \coordinate (A3) at (-1.75,-0.75);
    \coordinate (Ba) at ( 1.75, 0.75);
    \coordinate (Bb) at ( 1.75, 0);
    \coordinate (Bc) at ( 1.75,-0.75);
    \foreach \p in {A1,A2,A3,Ba,Bb,Bc} \node[mapelem] at (\p) {};
    \node[maplabL] at (A1) {$1$};
    \node[maplabL] at (A2) {$2$};
    \node[maplabL] at (A3) {$3$};
    \node[maplabR] at (Ba) {$a$};
    \node[maplabR] at (Bb) {$b$};
    \node[maplabR] at (Bc) {$c$};
    % 対応（1→c, 2→b, 3→a）
    \draw[maparrow] (A1) -- (Bc);
    \draw[maparrow] (A2) -- (Bb);
    \draw[maparrow] (A3) -- (Ba);
  \end{tikzpicture}
\end{example}

\begin{example}{学籍番号と学生}{学籍番号と学生}
  あるクラスの学生の集合を$A$，そのクラスで使われている学籍番号の集合を$B$とし，各学生$a \in A$にその学籍番号$f(a) \in B$を対応させる写像$f \colon A \to B$を考える．
  異なる学生が同じ学籍番号を持つことはないから$f$は単射であり，どの学籍番号もある学生に割り当てられているから$f$は全射である．よって$f$は全単射である．
\end{example}

\begin{example}{単射・全射・全単射}{単射・全射・全単射}
  \begin{enumerate}[(A)]
    \item
          \[
            f \colon [0,\infty) \ni x \mapsto x^2  \in \mathbb{R}
          \]
          は単射であるが全射でない．
    \item
          \[
            g \colon \mathbb{R} \ni x \mapsto x^3 -3x^2 -9 \in \mathbb{R}
          \]
          は全射であるが単射でない．
    \item
          \[
            h  \colon \mathbb{R} \ni x \mapsto  2 \sin x \in  \mathbb{R}
          \]
          は全射でも単射でもない．しかし，
          \[
            h'  \colon [-\pi/2 , \pi/2 ]  \ni x \mapsto  2 \sin  x  \in [-2,2]
          \]
          は全単射である．
  \end{enumerate}
  \end{example}
\phantomsection
\subsection{像・逆像の基本公式}

像と逆像は，\partref{集合}で学んだ集合の演算（$\cup$，$\cap$，$\setminus$）とどのように関わるのだろうか．まず，逆像については，集合の演算がすべてそのまま保たれる．

\phantomsection
\begin{prop}{逆像と集合演算}{逆像と集合演算}
  写像$f \colon A \to B$と$B_1, B_2 \subset B$に対して，次が成り立つ：
  \begin{align*}
    f^{-1}(B_1 \cup B_2) & = f^{-1}(B_1) \cup f^{-1}(B_2), \\
    f^{-1}(B_1 \cap B_2) & = f^{-1}(B_1) \cap f^{-1}(B_2), \\
    f^{-1}(B \setminus B_1) & = A \setminus f^{-1}(B_1)
  \end{align*}
\end{prop}

\begin{tleftbar}
  \begin{proof}
    2番目の等式を示す．$a \in A$に対して，
    \begin{align*}
      a \in f^{-1}(B_1 \cap B_2)
      & \lequiv f(a) \in B_1 \cap B_2 \\
      & \lequiv \bigl( f(a) \in B_1 \bigr) \land \bigl( f(a) \in B_2 \bigr) \\
      & \lequiv \bigl( a \in f^{-1}(B_1) \bigr) \land \bigl( a \in f^{-1}(B_2) \bigr) \\
      & \lequiv a \in f^{-1}(B_1) \cap f^{-1}(B_2)
    \end{align*}
    である．よって両者は等しい．残りの等式も，同様に逆像の定義に沿って書き換えれば示される．
  \end{proof}
\end{tleftbar}

一方，像については事情が異なる．和集合は保たれるが，共通部分については包含関係しか成り立たない．

\phantomsection
\begin{prop}{像と集合演算}{像と集合演算}
  写像$f \colon A \to B$と$A_1, A_2 \subset A$に対して，次が成り立つ：
  \begin{align*}
    f(A_1 \cup A_2) & = f(A_1) \cup f(A_2), \\
    f(A_1 \cap A_2) & \subset f(A_1) \cap f(A_2)
  \end{align*}
  ただし，2番目の包含は，一般には等号にならない．
\end{prop}

\begin{tleftbar}
  \begin{proof}
    2番目の包含を示す．$b \in f(A_1 \cap A_2)$とすると，ある$a \in A_1 \cap A_2$が存在して$b = f(a)$である．$a \in A_1$かつ$a \in A_2$なので，$b = f(a)$は$f(A_1)$にも$f(A_2)$にも属する．よって$b \in f(A_1) \cap f(A_2)$である．

    次に，等号が成り立たない例を挙げる．$f \colon \mathbb{R} \ni x \mapsto x^2 \in \mathbb{R}$，$A_1 = \{1\}$，$A_2 = \{-1\}$とすると，$A_1 \cap A_2 = \varnothing$より
    \[
      f(A_1 \cap A_2) = f(\varnothing) = \varnothing
    \]
    である一方，
    \[
      f(A_1) \cap f(A_2) = \{1\} \cap \{1\} = \{1\} \ne \varnothing
    \]
    である．なお，1番目の等式も，像の定義に沿って書き換えれば同様に示される．これが証明すべきことであった．
  \end{proof}
\end{tleftbar}

なぜ等号が崩れるのかというと，この$f$が異なる元$1$と$-1$を\textbf{同じ元}$1$にうつしてしまうからである．実際，$f$が単射であれば，この現象は起こらない．

\phantomsection
\begin{prop}{単射と像の共通部分}{単射と像の共通部分}
  写像$f \colon A \to B$が単射ならば，任意の$A_1, A_2 \subset A$に対して，
  \[
    f(A_1 \cap A_2) = f(A_1) \cap f(A_2)
  \]
  が成り立つ．
\end{prop}

\begin{tleftbar}
  \begin{proof}
    包含$\subset$は\prref{像と集合演算}で示したので，$\supset$を示す．$b \in f(A_1) \cap f(A_2)$とすると，ある$a_1 \in A_1$と$a_2 \in A_2$が存在して，
    \[
      b = f(a_1) = f(a_2)
    \]
    である．$f$は単射なので，$f(a_1) = f(a_2)$から$a_1 = a_2$が従う．よって$a_1 \in A_1 \cap A_2$であり，$b = f(a_1) \in f(A_1 \cap A_2)$である．これが証明すべきことであった．
  \end{proof}
\end{tleftbar}

この命題は，単射という性質がなぜ重要なのかを示すひとつの例である．逆像がつねに集合演算と相性がよいのに対して，像のふるまいは写像の性質（単射かどうか）に応じて変わるのである．

\phantomsection
\subsection{写像の合成とwell-defined}

\partref{用語の定義}では，「対象の表し方によらず値が定まること」をwell-definedとよんだ．写像の言葉が使えるようになったいま，この概念をもう少し正確に述べ直しておこう．準備として，写像の合成を定義する．

\phantomsection
\begin{definition}{合成写像}{合成写像}
  写像$h \colon X \to Y$と$g \colon Y \to Z$に対して，各$x \in X$に$g(h(x)) \in Z$を対応させる写像を$h$と$g$の\emph{合成写像}\index{ごうせいしゃぞう@合成写像}といい，$g \circ h \colon X \to Z$とかく：
  \[
    (g \circ h)(x) \coloneqq g(h(x)) \qquad (x \in X)
  \]
\end{definition}

「先に$h$，次に$g$」を施す写像が$g \circ h$であり，記号の順序が施す順序と逆になることに注意しよう．

 
また，高校数学で学んだ「合成関数」は，合成写像の特別な場合である．身近な合成写像の例をいくつか見ておこう．
 
\begin{example}{加法の合成}{加法の合成}
  「$2$を足す」写像と「$5$を足す」写像
  \[
    h \colon \mathbb{R} \ni x \mapsto x + 2 \in \mathbb{R}, \qquad
    g \colon \mathbb{R} \ni x \mapsto x + 5 \in \mathbb{R}
  \]
  を考える．合成写像$g \circ h$は，
  \[
    (g \circ h)(x) = g(h(x)) = g(x + 2) = (x + 2) + 5 = x + 7
  \]
  であり，たとえば$(g \circ h)(3) = (3 + 2) + 5 = 10$である．
  すなわち，$3 + 2 + 5$という計算は，$3$に写像$h$，$g$を順に施したものと見ることができる．
  この例では$(h \circ g)(x) = (x + 5) + 2 = x + 7$でもあり，$g \circ h = h \circ g$が成り立つ．
\end{example}
 
\begin{example}{数列の合成}{数列の合成}
  数列$a_1, a_2, a_3, \ldots$は，番号$n$に第$n$項$a_n$を対応させる写像
  \[
    a \colon \mathbb{N} \ni n \mapsto a_n \in \mathbb{R}
  \]
  とみなすことができる\footnote{本書では$\mathbb{N} = \{1, 2, 3, \ldots\}$とする．}．たとえば$a_n = n^2$とし，$k \colon \mathbb{N} \ni n \mapsto 2n \in \mathbb{N}$とすると，
  \[
    (a \circ k)(n) = a(k(n)) = a_{2n} = (2n)^2 = 4n^2
  \]
  である．すなわち，$a \circ k$は元の数列から偶数番目の項$a_2, a_4, a_6, \ldots$だけを取り出した数列$4, 16, 36, \ldots$である．
  一方，$g \colon \mathbb{R} \ni x \mapsto 2x + 1 \in \mathbb{R}$との合成$g \circ a$は，
  \[
    (g \circ a)(n) = g(a_n) = 2n^2 + 1
  \]
  であり，各項を$2x + 1$で変換した数列$3, 9, 19, \ldots$である．
\end{example}
 
\begin{example}{関数の合成}{関数の合成}
  $f \colon \mathbb{R} \ni x \mapsto \sin x \in \mathbb{R}$，$g \colon \mathbb{R} \ni x \mapsto x^2 \in \mathbb{R}$とすると，
  \[
    (g \circ f)(x) = (\sin x)^2 = \sin^2 x, \qquad
    (f \circ g)(x) = \sin(x^2)
  \]
  である．これらは異なる関数であり（たとえば$x = \sqrt{\pi}$で値が異なる），$g \circ f \ne f \circ g$である．
  \exref{加法の合成}とは対照的に，合成の順序を入れ替えると一般には別の写像になる．
\end{example}

さて，次のような状況を考える．写像$f \colon X \to Z$と全射$h \colon X \to Y$が与えられていて，$Y$の元$y$に対して「$h(x) = y$を満たす$x$をひとつ選び，その$f(x)$を$y$の行き先とする」という規則で，写像$g \colon Y \to Z$をつくりたい．$h$は全射だから，どの$y \in Y$に対しても$h(x) = y$を満たす$x \in X$は少なくともひとつ存在する．しかし，そのような$x$は一般にはひとつとは限らない．$h(x_1) = h(x_2) = y$を満たす$x_1 \ne x_2$があったとき，$f(x_1) \ne f(x_2)$であれば，$g(y)$の値は「どちらの$x$を選んだか」によって変わってしまい，$g$は写像の定義（\deref{写像}）の「ただひとつ対応させる」という条件を満たさない．この状況を整理したものが，次の定義である．

\phantomsection
\begin{definition}[
    sidebyside,
    lefthand width=8.6cm,
    righthand width=4.6cm,
    sidebyside gap=4mm,
  ]{well-defined}{well-defined}
  写像$f \colon X \to Z$と全射$h \colon X \to Y$を考える．任意の$x_1, x_2 \in X$に対して，
  \[
    h(x_1) = h(x_2) \limp f(x_1) = f(x_2)
  \]
  が成り立つとき，すなわち「$h$で同じ元にうつる$x$たちは，$f$でも同じ元にうつる」とき，$f(x)$の値は\emph{$y = h(x)$を満たす$x \in X$のとり方によらない}という．

  このとき，各$y \in Y$に対して，$h(x) = y$を満たす$x$をひとつ選んで
  \[
    g(y) \coloneqq f(x)
  \]
  と定めると，$g \colon Y \to Z$は写像であり，$f = g \circ h$を満たす．この$g$は\emph{well-defined}\index{welldefined@well-defined}（よく定義されている）であるという．

  \tcblower

  \centering
  \begin{tikzpicture}[mapfig, >=Latex]
    \node (X) at (0, 0)    {$X$};
    \node (Y) at (0, -2.4) {$Y$};
    \node (Z) at (2.6, -1.2) {$Z$};
    \draw[->, draw=black!70, line width=0.55pt, shorten <=3pt, shorten >=3pt]
      (X) -- (Y) node[midway, left=1pt] {$h$};
    \draw[->, draw=black!70, line width=0.55pt, shorten <=3pt, shorten >=3pt]
      (X) -- (Z) node[midway, above right=-1pt] {$f$};
    \draw[->, dashed, draw=black!70, line width=0.55pt, shorten <=3pt, shorten >=3pt]
      (Y) -- (Z) node[midway, below right=-1pt] {$g$};
  \end{tikzpicture}
  \par\smallskip
  {\small $f = g \circ h$}
\end{definition}

右の図は\emph{可換図式}\index{かかんずしき@可換図式}とよばれるもので，「$X$から$Z$へ，どの矢印をたどって行っても結果が同じ（$f = g \circ h$）」ことを表している．破線の$g$が「これから定めたい写像」であり，well-definedであるとは，この破線を実線に置き換えてよい，ということにほかならない．

\begin{remark}[「表し方によらない」との関係]
  \partref{用語の定義}で述べた「対象の表し方によらず値が定まる」という説明と，\deref{well-defined}とのつながりを確認しよう．$X$を「表し方」の集合，$Y$を「本来の対象」の集合とし，$h \colon X \to Y$を「表し方から対象を読み取る」写像だと思えばよい．同じ対象$y$は，$h(x) = y$を満たす$x$の個数だけ表し方をもつ．$f$は表し方$x$に対して定められた計算であり，
  「$f(x)$が表し方$x$のとり方によらない」ならば，$f$は本来の対象$y$に対する写像$g$を定めている\,---\,これが\deref{well-defined}のいうことである．
\end{remark}

\begin{example}{有理数の関数のwell-defined性}{有理数の関数のwell-defined性}
  $X = \{ (a, b) \in \mathbb{Z} \times \mathbb{Z} \mid b \ne 0 \}$を「分数の表し方」の集合，$Y = \mathbb{Q}$とし，
  \[
    h \colon X \ni (a, b) \mapsto a/b \in \mathbb{Q}
  \]
  とする．どの有理数も分数の形に表せるから，$h$は全射である．
  また，$h(1, 2) = h(2, 4)$のように，同じ有理数にうつる組は複数ある．
  \begin{enumerate}[(A)]
    \item $f \colon X \ni (a, b) \mapsto a + b \in \mathbb{Z}$とする．$h(1, 2) = h(2, 4)$であるが，$f(1, 2) = 3 \ne 6 = f(2, 4)$である．よって$f(a, b)$は$(a, b)$のとり方によるから，$g(a/b) = a + b$は写像$g \colon \mathbb{Q} \to \mathbb{Z}$を定めない（well-definedでない）．
    \item $f \colon X \ni (a, b) \mapsto 2a/b \in \mathbb{Q}$とする．$h(a, b) = h(a', b')$，すなわち$a/b = a'/b'$ならば，$ab' = a'b$より$2a/b = 2a'/b'$，すなわち$f(a, b) = f(a', b')$である．よって$g (a/b ) = 2a/b$はwell-definedであり，写像$g \colon \mathbb{Q} \to \mathbb{Q}$を定める．
  \end{enumerate}
  \partref{用語の定義}で述べた例が，まさにこの形をしていたことを確認してほしい．
\end{example}

\begin{remark}[「細かい」という言い回し]
  \deref{well-defined}の条件「$h(x_1) = h(x_2)$ならば$f(x_1) = f(x_2)$」は，$X$を「$h$の値が等しい元どうし」で分類したとき，その分類が「$f$の値が等しい元どうし」の分類より\emph{細かい}，と言い表されることがある．同じ$h$の値をもつ元たちは，かならず同じ$f$の値をもつ（が，逆は要求しない）からである．同値関係や商集合の言葉を学ぶと，この定義はさらに簡潔に述べ直すことができる．
\end{remark}

% NOTE: この部は執筆中（今後，節を追記する可能性あり）
